\section{Foundations: What Is a ``Causal'' Embedding?}
\label{sec:foundations}

The phrase ``causal embedding'' is used loosely across the retrieval and
representation-learning literature, often as a synonym for any embedding whose
geometry encodes \emph{asymmetric} relations. This section pins down what we
mean precisely, separates three distinct notions that are routinely conflated,
and locates the project's goal --- directional precedence between documentation
steps --- on the firmest defensible footing. We are deliberately conservative:
we argue that a ``causal embedding'' for procedural documentation is, in the
honest analysis, a model of \emph{structural / directional precedence learned
from text}, sitting almost entirely at the lowest rung of Pearl's ladder, and we
say so plainly rather than overclaiming interventional or counterfactual
content.

\subsection{Three notions routinely conflated}
\label{sec:three-notions}

\paragraph{(i) Statistical / semantic similarity.}
The default embedding objective --- contrastive or masked-language pretraining,
then a cosine/dot-product similarity --- produces a \emph{symmetric},
\emph{undirected} notion of relatedness: $\mathrm{sim}(a,b)=\mathrm{sim}(b,a)$.
Two documentation steps that mention the same dataset, schema, or UI panel land
near each other regardless of which must happen first. This is exactly the
shortcut our project must defeat: in the repo's hardest ordering dataset
(\texttt{ajo\_newstyle}), the best symmetric similarity model collapses to
$\approx 65\%$ accuracy because topical overlap and step order are decorrelated
on purpose. Semantic similarity is a \emph{necessary} substrate (it gives a
strong stage-1 retriever) but is provably insufficient for ordering: any
symmetric kernel assigns the pair $(A,B)$ and $(B,A)$ identical scores and
therefore carries zero bits about direction.

\paragraph{(ii) Asymmetric / directional relations.}
The next notion drops symmetry without yet claiming causality. Entailment is the
canonical example: ``$P$ entails $H$'' does not imply ``$H$ entails $P$''
\citep{found:bowman2015snli,found:williams2018mnli}. Order embeddings,
hyperbolic/entailment cones, and asymmetric dual encoders (a query tower and a
document tower with untied weights) all realize a scoring function
$s(a\!\to\!b)\neq s(b\!\to\!a)$. Temporal precedence (``$A$ occurs before $B$'')
and prerequisite structure also live here. Crucially, \emph{asymmetry is not
causality}: an asymmetric relation can encode narration order, set inclusion, or
specificity, none of which is a causal mechanism. Most so-called ``causal
embeddings'' in practice are this --- a learned directional relation --- and the
honest claim stops here.

\paragraph{(iii) Genuine causal structure.}
The strongest notion requires that the represented relation track an underlying
data-generating mechanism: $A\to B$ means intervening on $A$ changes the
distribution of $B$, while intervening on $B$ does not change $A$
\citep{found:pearl2009causality,found:peters2017elements}. This is defined with
respect to a structural causal model (SCM) and is identified, in the
observational regime, only under assumptions (Sec.~\ref{sec:direction-id}). For
text, we essentially never observe interventions on documentation steps, so
genuine causal structure is at best \emph{inferred} from the asymmetries the
text exposes. Our position: the project targets notion~(ii) as a tractable proxy
for notion~(iii), and the contribution is making that proxy as faithful as the
text allows.

\paragraph{Problem statement.}
Formally, given an ordered pair of texts $(t_1,t_2)$ we want a model
$f_\theta(t_1,t_2)\to[0,1]$ estimating $P(t_1 \text{ causally precedes /
enables } t_2)$, with the constraints (a) \emph{antisymmetry},
$f_\theta(t_1,t_2)\approx 1-f_\theta(t_2,t_1)$ for genuinely directional pairs,
and (b) an \emph{abstain / NEUTRAL} option for independent or parallel steps so
that the induced relation is a partial order (DAG), not a forced total order.
Downstream we need this primitive to (1) classify direction, (2) rerank
follow-ups, and (3) sort a workflow. A pure bi-encoder embedding cannot satisfy
(a) by construction unless direction is injected via untied towers or an
asymmetric scoring head; this motivates the cross-encoder / interaction-based
designs surveyed later.

\subsection{Causal representation learning: what it gives and does not give for text}
\label{sec:crl}

The modern program of \emph{causal representation learning} (CRL)
\citep{found:scholkopf2021causalrl} seeks to recover latent variables that
correspond to causal factors of variation and to learn the causal graph among
them, rather than entangled statistical features. Three pillars are relevant.

\paragraph{Independent causal mechanisms (ICM).}
The ICM principle holds that the causal generative process factorizes into
autonomous modules $P(X)=\prod_i P(X_i\mid \mathrm{pa}_i)$ whose mechanisms are
independent: changing one mechanism (an intervention) leaves the others
invariant \citep{found:peters2017elements,found:scholkopf2021causalrl}. This
\emph{invariance} is the deepest reason causal features generalize across
distribution shift, and it underlies invariant-risk and domain-generalization
objectives \citep{found:arjovsky2019irm,found:gulrajani2021domainbed}. For our
setting, ICM is aspirational: it suggests that a step's \emph{enabling
preconditions} should be an invariant module reusable across workflows, which is
precisely the generalization we need for far-apart (non-adjacent) step pairs.

\paragraph{Disentanglement and identifiability.}
A central negative result is that disentangled representations are
\emph{not identifiable} from i.i.d.\ observations without inductive biases or
weak supervision \citep{found:locatello2019challenging}: infinitely many latent
factorizations explain the same marginal. CRL therefore leans on auxiliary
signals --- multiple environments, interventions, temporal structure, or paired
data --- to pin down the causal factors \citep{found:scholkopf2021causalrl}.
\textbf{What this means for text:} we have no controlled interventions on
documentation, but we \emph{do} have weak supervision (the antisymmetric
forward/reverse labels, multiple documents, tier metadata). The identifiability
literature warns that direction learned from a single narration order is
under-determined --- it can latch onto narration rather than mechanism --- which
is exactly the failure diagnosed in the repo data (training only on adjacent,
same-document pairs teaches narration order, not causal necessity). CRL thus
does \emph{not} hand us a causal text embedding for free; it tells us which
auxiliary signals (non-adjacent pairs, multiple sources, abstention) are needed
to make direction identifiable.

\subsection{Pearl's ladder and where a documentation ``causal embedding'' sits}
\label{sec:ladder}

Pearl's \emph{ladder of causation} stratifies causal questions into three rungs
\citep{found:pearl2009causality,found:pearl2018book}:

\begin{description}[leftmargin=1.2em,itemsep=2pt]
  \item[Rung 1 --- Association] $P(y\mid x)$: ``seeing.'' Correlation,
  conditional probability, similarity. All standard embeddings live here.
  \item[Rung 2 --- Intervention] $P(y\mid \mathrm{do}(x))$: ``doing.'' The effect
  of actively setting $x$, requiring an SCM or experiments.
  \item[Rung 3 --- Counterfactual] $P(y_x\mid x',y')$: ``imagining.'' What would
  have happened had $x$ been different, holding the realized world fixed.
\end{description}

\noindent\textbf{An honest placement.}
A causal embedding for procedural documentation operates predominantly at
\emph{Rung 1 with structural / directional precedence layered on top}. We are
estimating an associational quantity --- $P(\text{$t_1$ precedes $t_2$}\mid
t_1,t_2)$ --- that we \emph{interpret} as enablement. We are not running
$\mathrm{do}(\cdot)$ on a workflow, and we are not computing counterfactual
re-orderings of a tutorial. The directional supervision (forward/reverse labels)
nudges the representation toward the asymmetry that a Rung-2 mechanism would
induce, but the trained object remains a Rung-1 estimator of a directed relation.
This is the single most important caveat in the survey: claiming Rung-2/3
content for a text embedding trained on observational documentation would be an
overclaim. Where genuine intervention-style evidence exists in our data, it is
implicit (``you cannot publish before you authenticate'' is a logical/temporal
precondition recoverable from text, not an experiment).

\subsection{Causal direction identification (background, then pivot to text)}
\label{sec:direction-id}

Outside NLP, the \emph{bivariate} causal-direction problem --- given samples of
$(X,Y)$, decide $X\to Y$ vs.\ $Y\to X$ --- is the cleanest analog of our pairwise
task. Because the joint $P(X,Y)$ is symmetric, direction is identifiable only
under structural assumptions that break the symmetry:

\begin{itemize}[leftmargin=1.2em,itemsep=2pt]
  \item \textbf{Additive Noise Models (ANM)} \citep{found:hoyer2008anm}:
  assume $Y=f(X)+N$, $N\!\perp\!X$. A nonlinear $f$ typically admits an
  independent-noise fit in only one direction; one regresses both ways and tests
  residual independence (HSIC) to pick the direction with independent residuals.
  \item \textbf{LiNGAM} \citep{found:shimizu2006lingam}: linear non-Gaussian
  acyclic models identify direction via independent-component analysis on the
  non-Gaussian noise.
  \item \textbf{IGCI} \citep{found:janzing2012igci}: information-geometric causal
  inference assumes independence between the cause distribution and the
  mechanism, scoring an asymmetry in their ``alignment.''
\end{itemize}

\noindent The standard testbed is the \emph{T\"ubingen cause-effect pairs}
benchmark of $\approx 100$ expert-labeled real bivariate pairs
\citep{found:mooij2016distinguishing}; reported accuracies cluster around
$\sim$60--72\% for IGCI/ANM-style methods --- i.e.\ direction is hard even with
clean numeric data and explicit modeling assumptions. \textbf{The pivot to
text:} documentation gives us neither i.i.d.\ numeric samples nor a clean noise
model. Instead, the asymmetry-revealing signal is \emph{linguistic} ---
discourse connectives, temporal markers, prerequisite phrasing, and world
knowledge --- which is why the rest of this survey treats causal direction in
text as a supervised pairwise classification / entailment problem
(Sec.~\ref{sec:causal-reasoning}) rather than an ANM/IGCI estimation problem.
The bivariate literature contributes the conceptual lesson, not the algorithm:
direction must come from an assumption or supervision that breaks the symmetry of
$P(X,Y)$, and for text that assumption is encoded in the antisymmetric labels and
the cross-attention architecture.

\subsection{A taxonomy of ``causal'' embeddings}
\label{sec:taxonomy}

Table~\ref{tab:found-taxonomy} organizes the design space along two axes: the
\emph{relation type} the model can express (symmetric, directional, or
mechanism-grounded) and the \emph{architecture} that realizes it. The taxonomy
makes explicit that genuine causal grounding is not achievable by geometry alone
in the observational text regime; the best a feasible model does is a
faithfully directional, abstention-capable estimator (row 2--3), built on a
strong symmetric retriever (row 1).

\begin{table*}[t]
\centering
\footnotesize
\caption{Taxonomy of embedding/relation models for procedural documentation.
``Pearl rung'' is the highest rung the model honestly addresses in the
observational text regime. Direction can be \emph{injected} architecturally
(untied towers, cross-attention) or via the training signal (antisymmetric
labels).}
\label{tab:found-taxonomy}
\begin{tabular}{@{}p{2.6cm}p{3.0cm}p{2.2cm}p{4.6cm}p{2.4cm}@{}}
\toprule
\textbf{Relation type} & \textbf{Typical objective} & \textbf{Symmetry} &
\textbf{Architecture / realization} & \textbf{Pearl rung} \\
\midrule
Semantic similarity & Contrastive (InfoNCE), MLM & Symmetric &
Single-tower / tied bi-encoder; cosine score & Rung 1 (assoc.) \\
\addlinespace
Directional (learned) & Pairwise BCE on forward/reverse labels; entailment &
Antisymmetric & Untied dual encoder; order/cone embeddings; cross-encoder over
$[t_1;t_2]$ & Rung 1 + directional precedence \\
\addlinespace
Directional + abstain & As above $+$ NEUTRAL class; partial-order/DAG decoding &
Antisymmetric, partial & Cross-encoder $+$ 3-way head; rank aggregation
(Bradley--Terry, min-feedback-arc-set) & Rung 1 + structural precedence \\
\addlinespace
Mechanism-grounded (aspirational) & Invariance / ICM; counterfactual
consistency & Asymmetric, invariant & CRL latent SCM; reasoning trace with
grounded premises (VeriCoT-style) & Rungs 2--3 (\emph{not} attained from text
alone) \\
\bottomrule
\end{tabular}
\end{table*}

\paragraph{Relevance to causal embedding for AEP/AJO.}
This foundation fixes the target and its honest ceiling. The AEP/AJO project
needs row~2--3 of Table~\ref{tab:found-taxonomy}: a \emph{directional,
abstention-capable} precedence estimator built on a strong symmetric retriever
(row~1, where fine-tuned BGE-large already gives the best MRR in the repo). Three
foundational lessons carry forward. First, symmetric similarity cannot order
steps --- direction must be \emph{injected}, motivating cross-attention
backbones over independent embeddings. Second, the identifiability results
\citep{found:locatello2019challenging,found:scholkopf2021causalrl} explain
\emph{why} training on adjacent same-document pairs fails to generalize:
direction learned from a single narration order is under-determined, so we must
supply the missing auxiliary signal (non-adjacent transitive pairs, multiple
sources, an abstain class). Third, we should advertise the model as a Rung-1
structural-precedence estimator, not an interventional one --- the ICM/invariance
ideal (row~4) is the north star for generalization, but genuine Rung-2/3 content
is not recoverable from observational documentation. The next section turns to
how direction is actually learned and evaluated from text.
